% 86-39(인쇄 86쪽) — 개념 25 지시문의 확률밀도함수 y=f(x) 그래프: (-1,0)에서 (0,2/3)까지 오르고 (1,2/3)까지 수평(오른쪽 끝 안내 파선).
% 스캔 화질이 낮아 TikZ 로 다시 그림. 라벨은 원문에 있는 것만(2/3, -1, O, 1, x, y, y=f(x)) · 답(확률)은 그림에 넣지 않는다.
\begin{tikzpicture}[x=1.05cm, y=1.15cm, line width=0.6pt]
  \draw[->, >=stealth, line width=0.7pt] (-1.75,0) -- (1.95,0) node[below=1pt, font=\small] {$x$};
  \draw[->, >=stealth, line width=0.7pt] (0,-0.45) -- (0,1.25) node[left=1pt, font=\small] {$y$};
  \draw[densely dotted, line width=0.4pt] (1,0) -- (1,0.667);
  \draw (-1,0) -- (0,0.667) -- (1,0.667);
  \node[anchor=east, inner sep=1pt, font=\small] (twothird) at (-0.62,0.72) {$\dfrac{2}{3}$};
  \draw[->, >=stealth, line width=0.35pt] (twothird.east) to[bend left=22] (-0.05,0.667);
  \node[below=1pt, font=\small] at (-1,0) {$-1$};
  \node[below=1pt, font=\small] at (1,0) {$1$};
  \node[below left=-1pt and -1pt, font=\small] at (0,0) {$\mathrm{O}$};
  \node[anchor=west, font=\small] at (0.18,0.9) {$y=f(x)$};
\end{tikzpicture}
